\subsection{Scaling Analysis of the Projective Gap}
  \label{subsec:scaling_projective_gap}

  The inevitability of saturation can be established through a simple
  scaling argument.
  Let the update demand scale as
  \begin{equation}
    D(R) \propto R^{-3/2},
  \end{equation}
  as dictated by the free-fall timescale.

  The projective capacity is limited by the maximal admissible throughput
  under the invariant bound~$c_{\mathrm{BI}}$.
  Denote by $C(R)$ the total admissible update capacity within a region
  of size~$R$.
  Without assuming a specific geometric postulate, we parametrize
  \begin{equation}
    C(R) \propto R^{n},
  \end{equation}
  where $n \ge 0$ encodes the effective growth of independent relational
  channels permitted by the mapping~$\Pi$ at the epoch considered.

  The projective stress ratio behaves as
  \begin{equation}
    \Xi(R) = \frac{D(R)}{C(R)} \propto R^{-(3/2+n)}.
  \end{equation}
  For any physically reasonable $n \ge 0$, the ratio diverges as
  $R \rightarrow 0$.
  This guarantees the existence of a finite contraction scale
  $R_{\mathrm sat}$ at which $\Xi \sim 1$.
  Beyond this point, the demanded update rate exceeds the admissible
  projective throughput.

  The result is robust with respect to the precise value of~$n$.
  Even if capacity scaled volumetrically ($n=3$), the divergence would
  remain unavoidable.
  Saturation during sufficiently deep collapse is therefore a structural
  feature of the framework rather than a model-dependent artifact.
